\section*{Bab \ref{ch:b1:seq} --- Barisan}

\begin{solution}{exo:b1:seq:1}
$\Bigl|\dfrac{2n+1}{n+3} - 2\Bigr| = \dfrac{5}{n+3}$. Diberikan
$\varepsilon > 0$, ambillah $N > \frac 5\varepsilon - 3$ (menurut Archimedes):
maka untuk $n \geq N$, $\frac{5}{n+3} \leq \varepsilon$. Jadi limitnya
$2$.

$((-1)^n)$: \omterm{def:b1:seq:subsequence}{subbarisannya} $(u_{2n}) = (1)$ dan $(u_{2n+1}) =
(-1)$ konvergen ke limit yang berbeda, sehingga barisannya divergen
(\cref{prop:b1:seq:subsequences}). (Secara langsung: sebarang calon $\ell$
gagal untuk $\varepsilon = \frac12$, karena suku yang berurutan berjarak
$2$.)
\end{solution}

\begin{solution}{exo:b1:seq:2}
Dengan \omterm{def:b1:arith:divides}{membaginya} dengan $n^2$: $\dfrac{1 - 3/n + 1/n^2}{2 + 5/n^2} \to
\dfrac12$.

$\sqrt{n+1} - \sqrt n = \dfrac{1}{\sqrt{n+1} + \sqrt n} \to 0$
(lewat sekawannya).

$\dfrac{2^n + n^3}{3^n - n^2} = \dfrac{(2/3)^n + n^3/3^n}{1 -
n^2/3^n} \to \dfrac{0 + 0}{1 - 0} = 0$, dengan memakai $q^n \to 0$ untuk $\abs q
< 1$ beserta perbandingan \omterm{def:b1:poly:def}{polinomial} lawan geometri
(\cref{prop:b1:functions:powerrules}).

$5^{1/n} = \eu^{\frac{\ln 5}{n}} \to \eu^0 = 1$.
\end{solution}

\begin{solution}{exo:b1:seq:3}
Bernoulli: $(1+h)^n \geq 1 + nh$ untuk $h \geq -1$, lewat induksi ---
$(1+h)^{n+1} = (1+h)^n(1+h) \geq (1+nh)(1+h) = 1 + (n+1)h + nh^2
\geq 1 + (n+1)h$.

Untuk $0 < \abs q < 1$: tulislah $\frac{1}{\abs q} = 1 + h$, $h > 0$; maka
$\abs{q}^n = \frac{1}{(1+h)^n} \leq \frac{1}{1 + nh} \to 0$, dan
apitannya memberikan $q^n \to 0$ (adapun kasus $q = 0$ sepele). Untuk $q = 1$:
barisannya konstan, dengan limit $1$. Untuk $q = -1$: divergen
(\cref{exo:b1:seq:1}). Untuk $\abs q > 1$: $\abs q^n = (1 + h)^n \geq
1 + nh \to +\infty$, sehingga $(q^n)$ tak terbatas, jadi divergen (ke
$+\infty$ bila $q > 1$; sedangkan dengan tanda berselang-seling, tanpa limit, bila $q <
-1$).
\end{solution}

\begin{solution}{exo:b1:seq:4}
Titik tetapnya: $\ell = \frac{\ell + 3}{2}$ memberikan $\ell = 3$. Lalu
\[
v_{n+1} = u_{n+1} - 3 = \frac{u_n + 3}{2} - 3 = \frac{u_n - 3}{2} =
\frac{v_n}{2}:
\]
$(v_n)$ bersifat geometri dengan rasio $\frac12$, dan $v_0 = -3$. Jadi $u_n = 3 -
\frac{3}{2^n} \to 3$.
\end{solution}

\begin{solution}{exo:b1:seq:5}
$H_{2n} - H_n = \sum_{k=n+1}^{2n} \frac 1k \geq n \cdot
\frac{1}{2n} = \frac12$ (karena masing-masing dari $n$ sukunya bernilai $\geq
\frac{1}{2n}$). Seandainya $(H_n)$ Cauchy, mengambil $\varepsilon =
\frac13$ akan memaksa $\abs{H_{2n} - H_n} \leq \frac13$ untuk $n$ yang
besar: yang bertentangan. Adapun barisan naik yang tak konvergen divergen ke
$+\infty$ (\cref{thm:b1:seq:monotone}): jadi $H_n \to +\infty$.
\end{solution}

\begin{solution}{exo:b1:seq:6}
Misalkan $a = \lim u_{2n}$, $b = \lim u_{2n+1}$, $c = \lim u_{3n}$. Barisan
$(u_{6n})$ merupakan \omterm{def:b1:seq:subsequence}{subbarisan} $(u_{2n})$ sekaligus
$(u_{3n})$: sehingga limitnya sama dengan $a$ dan $c$, jadi $a = c$. Barisan
$(u_{6n+3})$ merupakan \omterm{def:b1:seq:subsequence}{subbarisan} $(u_{2n+1})$ (yaitu indeks ganjil) dan
$(u_{3n})$ (dengan indeks $6n + 3 = 3(2n+1)$): sehingga $b = c$. Jadi $a = b$,
lalu \cref{prop:b1:seq:subsequences} (yaitu genap dan ganjil dengan limit yang
sama) memberikan kekonvergenan $(u_n)$.
\end{solution}

\begin{solution}{exo:b1:seq:7}
\emph{Kestabilan dan batasnya:} $I = \intcc{0}{2}$ stabil: karena untuk $x \in
I$, $\sqrt{2 + x} \in \intcc{\sqrt 2}{2} \subseteq I$; dan $u_0 = 0
\in I$.

\emph{Kemonotonannya:} $f(x) = \sqrt{2+x}$ naik dan $u_1 =
\sqrt 2 > u_0$: sehingga lewat induksi $(u_n)$ naik. Naik dan
terbatas di atas oleh $2$: jadi ia konvergen
(\cref{thm:b1:seq:monotone}).

\emph{Limitnya:} $\ell = \sqrt{2 + \ell}$ dengan $\ell \geq 0$ memberikan
$\ell^2 - \ell - 2 = 0$, sehingga $\ell = 2$.

\emph{Batas galatnya:} dengan mengalikannya dengan sekawannya,
\[
2 - u_{n+1} = 2 - \sqrt{2 + u_n}
= \frac{4 - (2 + u_n)}{2 + \sqrt{2+u_n}}
= \frac{2 - u_n}{2 + \sqrt{2 + u_n}}
\leq \frac{2 - u_n}{3},
\]
karena $\sqrt{2 + u_n} \geq \sqrt 2 > 1$. Lalu lewat induksi dari $2 - u_0
= 2$: $\;0 \leq 2 - u_n \leq \frac{2}{3^n}$.
\end{solution}

\begin{solution}{exo:b1:seq:8}
Andaikan $(u_n)$ tak konvergen ke $\ell$: maka untuk suatu $\varepsilon_0
> 0$, tak hingga banyak indeks memenuhi $\abs{u_n - \ell} >
\varepsilon_0$; dan indeks itu membentuk \omterm{def:b1:seq:subsequence}{subbarisan} $(u_{\varphi(n)})$. \omterm{def:b1:seq:subsequence}{Subbarisan}
ini terbatas, sehingga menurut Bolzano--Weierstrass
(\cref{thm:b1:seq:bw}) ia mempunyai subsubbarisan yang konvergen, yang
limitnya $\ell'$ memenuhi $\abs{\ell' - \ell} \geq \varepsilon_0$
(karena ketaksamaannya diteruskan ke limitnya, \cref{thm:b1:seq:order}). Padahal
subsubbarisan $(u_n)$ merupakan \omterm{def:b1:seq:subsequence}{subbarisan} $(u_n)$ yang konvergen,
sehingga menurut hipotesisnya $\ell' = \ell$: yang bertentangan.
\end{solution}

\begin{solution}{exo:b1:seq:9}
Andaikan $\eu = \frac pq$, dengan $q \geq 1$. Ketaksamaan tegas $a_q <
\eu < a_q + \frac{1}{q\,q!}$ (yang tegas karena $(a_n)$ naik tegas
dan $(b_n)$ turun tegas) bila dikalikan dengan $q!$ memberikan
\[
q!\,a_q \;<\; q!\,\frac pq \;<\; q!\,a_q + \frac 1q \leq q!\,a_q + 1.
\]
Sekarang $N = q!\,a_q = \sum_{k=0}^{q} \frac{q!}{k!}$ merupakan bilangan bulat (karena setiap
$\frac{q!}{k!}$ adalah hasil kali bilangan bulat untuk $k \leq q$), dan demikian pula
$q!\,\frac pq = (q-1)!\,p$. Jadi tampilannya menaruh bilangan bulat
$(q-1)!\,p$ tegas di antara $N$ dan $N + \frac 1q \leq N + 1$: yaitu sebuah
bilangan bulat yang tegas berada di dalam $\intoo{N}{N+1}$, yang mustahil. Jadi
$\eu \notin \Q$.
\end{solution}

\begin{solution}{exo:b1:seq:10}
\begin{enumerate}
  \item Misalkan $\varepsilon > 0$ dan $N$ dengan $\abs{u_k - \ell} \leq
        \frac{\varepsilon}{2}$ untuk $k > N$. Untuk $n > N$:
        \[
        \abs{c_n - \ell}
        = \Bigl|\frac{\sum_{k=1}^{n}(u_k - \ell)}{n}\Bigr|
        \leq \frac{\sum_{k=1}^{N} \abs{u_k - \ell}}{n}
        + \frac{n - N}{n}\cdot\frac{\varepsilon}{2}
        \leq \frac{C}{n} + \frac{\varepsilon}{2},
        \]
        dengan $C = \sum_{k=1}^N \abs{u_k - \ell}$ yang tetap. Untuk $n$
        yang besar, $\frac Cn \leq \frac\varepsilon2$: maka $\abs{c_n -
        \ell} \leq \varepsilon$.
  \item $u_n = (-1)^n$: divergen, namun $c_n \to 0$ (karena jumlah parsialnya
        terbatas oleh $1$, lalu dibagi $n$).
  \item Terapkan (1) pada barisan $v_n = u_{n+1} - u_n \to \ell$: maka
        rata-rata Cesàronya adalah $\frac{u_{n+1} - u_1}{n} \to \ell$
        (secara teleskopis), dan $\frac{u_{n+1}}{n} = \frac{u_{n+1} -
        u_1}{n} + \frac{u_1}{n} \to \ell$; lalu menormalkan ulang indeksnya
        ($\frac{u_n}{n} = \frac{u_n}{n-1}\cdot\frac{n-1}{n}$) memberikan
        $\frac{u_n}{n} \to \ell$.
\end{enumerate}
\end{solution}

\begin{solution}{exo:b1:seq:11}
Misalkan $L = \inf_{n \geq 1} \frac{u_n}{n} \geq 0$, dan $\varepsilon >
0$. Pilihlah $m$ dengan $\frac{u_m}{m} \leq L + \varepsilon$. Setiap $n$
tertulis $n = qm + r$, $0 \leq r < m$; lalu subaditifnya (yang diulang) memberikan
$u_n \leq q\,u_m + u_r$, sehingga
\[
\frac{u_n}{n} \leq \frac{q m}{n}\cdot\frac{u_m}{m} + \frac{u_r}{n}
\leq \frac{u_m}{m} + \frac{\max(u_0, \dots, u_{m-1})}{n}
\leq L + \varepsilon + \frac{C_m}{n},
\]
dengan memakai $qm \leq n$. Untuk $n$ yang besar, $\frac{C_m}{n} \leq \varepsilon$:
jadi $L \leq \frac{u_n}{n} \leq L + 2\varepsilon$ untuk setiap $n$ yang besar,
dan itulah kekonvergenan ke $L$.
\end{solution}

\begin{solution}{exo:b1:seq:12}
\omterm{def:b1:structures:subgroup}{Subgrup} $G = \Z + 2\pi\Z$ dari $(\R, +)$ bersifat padat: karena ia bukan
$\alpha\Z$, sebab $1 = p\alpha$, $2\pi = q\alpha$ akan membuat $2\pi =
\frac qp$ rasional --- padahal $\pi \notin \Q$ (yang diterima di sini; adapun buktinya
disketsakan pada \cref{ch:b1:integration}). Menurut \cref{exo:b1:reals:9}, $G$
padat di $\R$.

Sekarang misalkan $y \in \intcc{-1}{1}$ dan $\theta = \arcsin y$. Menurut kepadatannya,
untuk setiap $\varepsilon > 0$ ada $n \in \Z$, $k \in \Z$ dengan
$\abs{(n + 2\pi k) - \theta} \leq \varepsilon$, yakni $n$ berada dalam jarak
$\varepsilon$ dari $\theta - 2\pi k$; lalu, karena $\sin$ bersifat
periodik-$2\pi$ dan Lipschitz-$1$ ($\abs{\sin a - \sin b} \leq
\abs{a - b}$, dari ketaksamaan nilai rata-rata pada
\cref{ch:b1:derivative}),
\[
\abs{\sin n - y} = \abs{\sin(n + 2\pi k) - \sin\theta}
\leq \abs{n + 2\pi k - \theta} \leq \varepsilon .
\]
Satu perincian: $n$ menjelajahi $\Z$, tetapi $\sin(-n) = -\sin n$ dan $y$
tadi sebarang di $\intcc{-1}{1}$, sehingga indeks yang taknegatif sudah mencukupi
(gantilah $(n, y)$ dengan $(-n, -y)$ bila perlu). Jadi $\{\sin n : n \in
\N\}$ padat di $\intcc{-1}{1}$; dan barisan yang padat pada sebuah ruas mempunyai
\omterm{def:b1:seq:subsequence}{subbarisan} yang mendekati nilai yang berbeda, sehingga ia divergen.
\end{solution}

\begin{solution}{pb:b1:seq:1}
\textbf{1.} Di sini $\dfrac{n(n+1)/2}{n^2} = \dfrac{1 + 1/n}{2} \to
\dfrac12$, dan $\dfrac{n(n+1)(2n+1)/6}{n^3} = \dfrac{(1 +
1/n)(2 + 1/n)}{6} \to \dfrac13$.

\textbf{2.} Batas atasnya: masing-masing dari $n$ sukunya bernilai $\leq \sqrt n$, sehingga
$T_n \leq n\sqrt n$. Batas bawahnya: suku dengan $k > \frac n2$ berjumlah
sekurang-kurangnya $\frac n2$, dan masing-masingnya $\geq \sqrt{n/2}$:
\[
T_n \geq \frac n2 \sqrt{\frac n2} = \frac{n^{3/2}}{2\sqrt 2} .
\]

\textbf{3.} Untuk $k \geq N$, karena $b_{k+1} - b_k > 0$:
$m(b_{k+1} - b_k) \leq a_{k+1} - a_k \leq M(b_{k+1} - b_k)$.
Lalu menjumlahkannya untuk $k = N, \dots, n - 1$, kedua ruasnya teleskopis:
\[
m\,(b_n - b_N) \leq a_n - a_N \leq M\,(b_n - b_N),
\]
dan \omterm{def:b1:arith:divides}{membaginya} dengan $b_n - b_N > 0$ memberikan klaimnya.

\textbf{4.} Menurut definisi limitnya ada $N$ dengan $\ell -
\varepsilon \leq \frac{a_{k+1} - a_k}{b_{k+1} - b_k} \leq \ell +
\varepsilon$ untuk setiap $k \geq N$; lalu pertanyaan 3 dengan $m = \ell -
\varepsilon$, $M = \ell + \varepsilon$ memindahkan apitannya
ke $\frac{a_n - a_N}{b_n - b_N}$.

\textbf{5.} Dengan menjabarkan ruas kanan kesamaannya:
\[
\frac{a_N - \ell b_N}{b_n} + \frac{a_n - a_N}{b_n} -
\ell\,\frac{b_n - b_N}{b_n}
= \frac{a_n - \ell b_n}{b_n} = \frac{a_n}{b_n} - \ell .
\]
Menurut pertanyaan 4, faktor kedua hasil kalinya terbatas oleh
$\varepsilon$ dalam nilai mutlaknya, dan $0 < 1 - \frac{b_N}{b_n}
\leq 1$ untuk $n$ yang besar, sehingga
\[
\Bigl|\frac{a_n}{b_n} - \ell\Bigr|
\leq \frac{\abs{a_N - \ell b_N}}{b_n} + \varepsilon
\leq 2\varepsilon
\]
begitu $b_n \geq \frac{\abs{a_N - \ell b_N}}{\varepsilon}$,
yang akhirnya terjadi karena $b_n \to +\infty$. Jadi
$\frac{a_n}{b_n} \to \ell$: itulah teorema Cesàro--Stolz.

\textbf{6.} Diberikan $M$, pilihlah $N$ dengan $\frac{a_{k+1} -
a_k}{b_{k+1} - b_k} \geq M$ untuk $k \geq N$; lalu paruh bawah
pertanyaan 3 memberikan $a_n - a_N \geq M(b_n - b_N)$, sehingga
\[
\frac{a_n}{b_n} \geq \frac{a_N}{b_n} +
M\Bigl(1 - \frac{b_N}{b_n}\Bigr) \longrightarrow M .
\]
Di luar suatu peringkat, $\frac{a_n}{b_n} \geq M - 1$; dan karena $M$ tadi
sebarang, $\frac{a_n}{b_n} \to +\infty$.

\textbf{7.} Dengan $b_n = n$ dan $a_n = u_1 + \dots + u_n$: hasil bagi
selisihnya adalah $u_{n+1} \to \ell$, sehingga rata-rata Cesàro
$\frac{a_n}{n}$ menuju $\ell$: yaitu bagian (1)
\cref{exo:b1:seq:10}. Dengan $a_n = u_n$: hasil bagi selisihnya
adalah $u_{n+1} - u_n$, yang memberikan bagian (3). Konversnya: $a_n = (-1)^n$,
$b_n = n$ mempunyai $\frac{a_n}{b_n} \to 0$, namun $a_{n+1} - a_n = \pm
2$ berselang-seling: sehingga hasil bagi selisihnya tak berlimit.

\textbf{8.} Lewat pensekawanan:
\[
\bigl((1+h)^{3/2} - 1\bigr)\bigl((1+h)^{3/2} + 1\bigr)
= (1+h)^3 - 1 = 3h + 3h^2 + h^3 .
\]
Untuk $h = \frac1n$: $n\bigl((1 + \frac1n)^{3/2} - 1\bigr) =
\frac{3 + 3/n + 1/n^2}{(1 + 1/n)^{3/2} + 1}$, dan $1 \leq (1 +
\frac1n)^{3/2} \leq (1 + \frac1n)^2 \to 1$ (lewat apitan), sehingga
penyebutnya menuju $2$ dan ungkapannya menuju $\frac32$. Sekarang
Stolz dengan $a_n = T_n$, $b_n = n^{3/2}$ (yang naik tegas dan
$\to +\infty$):
\[
\frac{\sqrt{n+1}}{(n+1)^{3/2} - n^{3/2}}
= \frac{\sqrt{n+1}}{\sqrt n} \cdot
\frac{1}{n\bigl((1 + \frac1n)^{3/2} - 1\bigr)}
\longrightarrow 1 \cdot \frac{2}{3},
\]
sehingga $T_n \sim \frac23\,n^{3/2}$. (Adapun apitan pertanyaan 2
tadi memerangkap konstantanya di $\intcc{0.35}{1}$; sedangkan Stolz memakukannya.)

\textbf{9.} Dengan (G1) pada $u = \ln(1+t)$: $1 + t = \eu^{\ln(1+t)} \geq
1 + \ln(1+t)$, sehingga $\ln(1+t) \leq t$. Lalu (G1) pada $u =
-\frac{t}{1+t}$: $\eu^{-t/(1+t)} \geq 1 - \frac{t}{1+t} =
\frac{1}{1+t} > 0$; lalu mengambil $\ln$ (yang naik):
$-\frac{t}{1+t} \geq -\ln(1+t)$, yakni $\ln(1+t) \geq
\frac{t}{1+t}$.

\textbf{10.} Fungsi $\ln$ naik tegas
(\cref{prop:b1:functions:expln}), dan $\ln(2^k) = k\ln 2$
tak terbatas, sehingga $\ln n \to +\infty$. Selisihnya: dengan $t =
\frac1n$ pada pertanyaan 9,
\[
\frac 1{n+1} = \frac{1/n}{1 + 1/n} \leq
\ln\Bigl(1 + \frac1n\Bigr) \leq \frac1n
\quad\Longrightarrow\quad
\frac{n}{n+1} \leq \frac{1/(n+1)}{\ln(1 + 1/n)} \leq 1 ,
\]
sehingga hasil bagi selisihnya $\frac{H_{n+1} - H_n}{\ln(n+1) - \ln
n}$ menuju $1$; lalu Stolz memberikan $H_n \sim \ln n$.

\textbf{11.} Tetapkan $x_n = \frac{u_{n+1}}{u_n} \to L$ dan $t_n =
\frac{x_n}{L} - 1 \to 0$. Menurut pertanyaan 9: $\frac{t_n}{1 + t_n} \leq
\ln(1 + t_n) \leq t_n$, sehingga $\ln x_n - \ln L = \ln(1 + t_n) \to
0$ lewat pengapitan. Lalu Cesàro (pertanyaan 7) yang diterapkan pada $(\ln x_k)$:
\[
\frac1n \sum_{k=0}^{n-1} \ln x_k = \frac{\ln u_n - \ln u_0}{n}
\longrightarrow \ln L ,
\]
sehingga $\frac{\ln u_n}{n} \to \ln L$. Dengan $h_n = \frac{\ln u_n}{n}
- \ln L \to 0$: $u_n^{1/n} = L\,\eu^{h_n}$, dan (G1) mengapit
$1 + h_n \leq \eu^{h_n} \leq \frac{1}{1 - h_n}$ (untuk $h_n < 1$),
sehingga $\eu^{h_n} \to 1$ dan $u_n^{1/n} \to L$. Penerapannya: $u_n =
\binom{2n}{n}$ memberikan
\[
\frac{u_{n+1}}{u_n} = \frac{(2n+1)(2n+2)}{(n+1)^2}
= \frac{2(2n+1)}{n+1} \longrightarrow 4 ,
\qquad\text{sehingga}\qquad
\binom{2n}{n}^{1/n} \to 4 .
\]

\textbf{12.} Untuk $0 < x \leq 1$, (G2) memberikan $\sin x \geq x(1 -
\frac{x^2}{6}) \geq \frac{5x}{6} > 0$ dan
\[
x - \sin x \geq \frac{x^3}{6} - \frac{x^5}{120}
= x^3\Bigl(\frac16 - \frac{x^2}{120}\Bigr)
\geq \frac{19}{120}\,x^3 > 0 :
\]
jadi $0 < \sin x < x$ pada $\intoc{0}{1}$. Selalu berlaku $u_1 = \sin u_0
\in \intcc{-1}{1}$. Jika $u_1 = 0$, maka $u_n = 0$ untuk $n \geq 1$.
Jika $u_1 \in \intoc{0}{1}$: maka lewat induksi $0 < u_{n+1} = \sin u_n
< u_n \leq 1$, sehingga $(u_n)_{n\geq1}$ turun tegas dan
terbatas di bawah oleh $0$: jadi ia konvergen ke suatu $\ell \in
\intco{0}{1}$ (\cref{thm:b1:seq:monotone}). Adapun rumus hasil kali menjadi
jumlah beserta (G2) memberikan $\abs{\sin a - \sin b} = 2\abs{\cos
\frac{a+b}{2}}\,\abs{\sin\frac{a-b}{2}} \leq \abs{a - b}$, sehingga
$u_{n+1} = \sin u_n \to \sin \ell$: jadi $\ell = \sin\ell$. Jika $\ell
> 0$ maka $\sin\ell < \ell$: yang mustahil. Jadi $u_n \to 0$.

\textbf{13.} Dengan \omterm{def:b1:arith:divides}{membagi} (G2) dengan $u_n > 0$:
\[
1 - \frac{u_n^2}{6} \leq \frac{\sin u_n}{u_n} \leq 1 -
\frac{u_n^2}{6} + \frac{u_n^4}{120} \leq 1 ,
\]
dan $u_n \to 0$ mengapit $\frac{\sin u_n}{u_n} \to 1$. Lalu dengan membagi
$x - \sin x$ dengan $x^3$:
\[
\frac16 - \frac{u_n^2}{120} \leq \frac{u_n - \sin u_n}{u_n^3}
\leq \frac16 \longrightarrow \frac16 .
\]

\textbf{14.} Karena $u_{n+1} = \sin u_n$:
\[
w_n = \frac{u_n^2 - \sin^2 u_n}{u_n^2 \sin^2 u_n}
= \frac{(u_n - \sin u_n)(u_n + \sin u_n)}{u_n^2 \sin^2 u_n}
= \frac{u_n - \sin u_n}{u_n^3}\cdot
\frac{u_n + \sin u_n}{u_n}\cdot
\Bigl(\frac{u_n}{\sin u_n}\Bigr)^{2}
\]
(periksalah pangkat $u_n$-nya: $3 + 1 + (-4)$ terhadap $u_n^2$
di bawah dan $u_n^4$ di atas). Menurut pertanyaan 13, ketiga
faktornya menuju $\frac16$, $2$, $1$: jadi $w_n \to \frac13$.

\textbf{15.} Di sini $v_n = \frac{1}{u_n^2}$ berselisih $v_{n+1} -
v_n = w_n \to \frac13$, sehingga $\frac{v_n}{n} \to \frac13$ menurut
\cref{exo:b1:seq:10}~(3): jadi $n u_n^2 \to 3$. Lalu
\[
\abs{\sqrt n\,u_n - \sqrt 3}
= \frac{\abs{n u_n^2 - 3}}{\sqrt n\,u_n + \sqrt 3}
\leq \frac{\abs{n u_n^2 - 3}}{\sqrt 3} \longrightarrow 0 :
\]
jadi $\sqrt n\,u_n \to \sqrt 3$, yakni $u_n \sim \sqrt{3/n}$.

\textbf{16.} Karena $n u_n^2 \to 3$, akhirnya $2 \leq n u_n^2
\leq 4$, yakni $\sqrt{2/n} \leq u_n \leq 2/\sqrt n$. Jika $u_n
\leq 10^{-2}$ dengan $n$ pada rentang itu, maka $2/n \leq 10^{-4}$:
jadi $n \geq 20\,000$; dan $\sqrt{3/n} = 10^{-2}$ pada $n = 30\,000$.
Adapun \omterm{ex:b1:seq:heron}{metode Heron} mengkuadratkan galatnya pada setiap langkah --- sehingga cacah
angkanya berlipat dua --- karena pada titik tetapnya kemiringan yang relevan
bermodulus $< 1$ (memang iterasinya bersifat mengontraksi).
Di sini $\sin' 0 = \cos 0 = 1$: jadi titik tetapnya netral, tak ada
kontraksi geometri, dan peluruhannya dikuasai oleh
suku taklinear yang pertama $-\frac{x^3}{6}$, sehingga bersifat \omterm{def:b1:poly:def}{polinomial}. Satu
langkah Heron memperoleh ketelitian yang lebih banyak daripada sepuluh ribu langkah
sinusnya.

\textbf{17.} Untuk $u_0$ yang sebarang: $u_1 = \sin u_0 \in
\intcc{-1}{1}$. Jika $u_1 = 0$ maka barisannya lenyap mulai peringkat
$1$. Jika $u_1 > 0$, Bagian IV berlaku persis apa adanya. Jika $u_1 < 0$, tetapkan
$v_n = -u_n$: maka keganjilan $\sin$ memberikan $v_{n+1} = -\sin u_n =
\sin(-u_n) = \sin v_n$ dengan $v_1 \in \intoc{0}{1}$, sehingga $v_n \sim
\sqrt{3/n}$, yakni $u_n \sim -\sqrt{3/n}$. Pada semua kasusnya
$\abs{u_n} \sim \sqrt{3/n}$ (atau barisannya akhirnya $0$):
jadi jatuhnya bersifat semesta, dan hanya tandanya yang mengingat awalnya.

\textbf{18.} Pertama $\frac{u_{n+1}}{u_n} = 1 - \frac{u_n -
u_{n+1}}{u_n^2}\,u_n \to 1 - a \cdot 0 = 1$. Lalu
\[
\frac{1}{u_{n+1}} - \frac{1}{u_n}
= \frac{u_n - u_{n+1}}{u_n u_{n+1}}
= \frac{u_n - u_{n+1}}{u_n^2}\cdot\frac{u_n}{u_{n+1}}
\longrightarrow a \cdot 1 = a ,
\]
dan \cref{exo:b1:seq:10}~(3) memberikan $\frac{1}{n u_n} \to a$,
yakni $n u_n \to \frac1a$.

\textbf{19.} Dengan $v_n = \frac{1}{u_n}$: $v_{n+1} = \frac{1 +
u_n}{u_n} = v_n + 1$, sehingga $v_n = v_0 + n$ dan
\[
u_n = \frac{u_0}{1 + n u_0} , \qquad
n u_n = \frac{n u_0}{1 + n u_0} \longrightarrow 1 .
\]
Periksa lemanya: $u_n - u_{n+1} = \frac{u_n^2}{1 + u_n}$, sehingga
$\frac{u_n - u_{n+1}}{u_n^2} = \frac{1}{1 + u_n} \to 1 = a$, dan
pertanyaan 18 meramalkan $n u_n \to 1$: jadi persis bersesuaian.

\textbf{20.} Kepositifannya lewat induksi (karena $\eu^{-u} > 0$);
ia turun karena $\eu^{-u_n} < 1$ untuk $u_n > 0$; sehingga $u_n \to
\ell \geq 0$ (\cref{thm:b1:seq:monotone}). Jembatan kekontinuannya:
dengan $h_n = \ell - u_n \to 0$, $\eu^{-u_n} = \eu^{-\ell}
\eu^{h_n} \to \eu^{-\ell}$ menurut apitan (G1) $1 + h_n \leq
\eu^{h_n} \leq \frac{1}{1 - h_n}$; jadi $\ell = \ell\,
\eu^{-\ell}$, dan $\ell > 0$ akan memaksa $\eu^{-\ell} = 1$,
yang salah: sehingga $\ell = 0$. Untuk $t > 0$, (G1) memberikan $\eu^{-t} \geq 1 - t$
dan $\eu^{-t} \leq \frac{1}{1 + t}$, sehingga
\[
\frac{1}{1 + t} \leq \frac{1 - \eu^{-t}}{t} \leq 1 .
\]
Dengan $t = u_n$: $\frac{u_n - u_{n+1}}{u_n^2} = \frac{1 -
\eu^{-u_n}}{u_n} \to 1$. Lalu pertanyaan 18 dengan $a = 1$: $n u_n \to
1$, sehingga $u_n \sim \frac1n$.

\textbf{21.} Seperti pada pertanyaan 18, $\frac{u_{n+1}}{u_n} = 1 -
\frac{u_n - u_{n+1}}{u_n^3}\,u_n^2 \to 1$. Lalu
\[
\frac{1}{u_{n+1}^2} - \frac{1}{u_n^2}
= \frac{(u_n - u_{n+1})(u_n + u_{n+1})}{u_n^2 u_{n+1}^2}
= \frac{u_n - u_{n+1}}{u_n^3}\cdot
\frac{u_n + u_{n+1}}{u_n}\cdot
\Bigl(\frac{u_n}{u_{n+1}}\Bigr)^2
\longrightarrow a \cdot 2 \cdot 1 = 2a ,
\]
dan \cref{exo:b1:seq:10}~(3) memberikan $\frac{1}{n u_n^2} \to 2a$:
jadi $n u_n^2 \to \frac{1}{2a}$. Untuk sinusnya, $a = \frac16$
(pertanyaan 13): sehingga $n u_n^2 \to 3$, yaitu persis Bagian IV.

\textbf{22.} Selisihnya: $a_{n+1} - a_n = 2^{-n} - 2^{-n-1} =
2^{-n-1} = b_{n+1} - b_n$, sehingga hasil bagi selisihnya
selalu $1$. Namun $\frac{a_n}{b_n} = \frac{2 - 2^{-n}}{1 -
2^{-n}} \to 2 \neq 1$. Adapun bukti pertanyaan 5 patah pada
suku batasnya: karena $\frac{a_N - \ell b_N}{b_n} \to 0$ memerlukan $b_n
\to +\infty$; sedangkan di sini (dengan $\ell = 1$) $a_N - b_N = 1$ dan $b_n
\to 1$, sehingga sukunya menuju $1$ --- yaitu persis celah sisanya
$2 - 1$.

\textbf{23.} Stolz dengan $A_n = \sum_{k=1}^n H_k$ dan $B_n = n\ln
n$: di sini $B_{n+1} - B_n = \ln(n+1) + n\ln(1 + \frac1n) > 0$ dan $B_n
\to +\infty$. Menurut pertanyaan 9, $\frac{n}{n+1} \leq n\ln(1 +
\frac1n) \leq 1$, sehingga $B_{n+1} - B_n = \ln(n+1) + \theta_n$ dengan
$\frac12 \leq \theta_n \leq 1$. Jadi
\[
\frac{A_{n+1} - A_n}{B_{n+1} - B_n}
= \frac{H_{n+1}}{\ln(n+1)}\cdot
\frac{1}{1 + \theta_n/\ln(n+1)} \longrightarrow 1 \cdot 1 = 1
\]
(dengan pertanyaan 10 untuk faktor pertamanya; sedangkan $\theta_n$ terbatas dan
$\ln(n+1) \to \infty$ untuk yang kedua). Lalu Stolz menyimpulkan:
$\sum_{k=1}^n H_k \sim n\ln n$.

\textbf{24.} Jika $u_n \to \ell > 0$: maka seperti pada pertanyaan 11, $\ln u_n
\to \ln\ell$ (lewat apitan pertanyaan 9 pada $\ln\frac{u_n}{\ell}$), sehingga
rata-rata Cesàronya $\frac1n\sum_{k=1}^n \ln u_k \to \ln\ell$,
dan jembatan eksponensial pertanyaan 11 memberikan $(u_1\cdots
u_n)^{1/n} = \exp\bigl(\frac1n\sum\ln u_k\bigr) \to \ell$. Jika
$u_n \to +\infty$: maka untuk sebarang $M$, akhirnya $u_n \geq \eu^M$, sehingga
$\ln u_n \geq M$: jadi $\ln u_n \to +\infty$; lalu Cesàro yang $+\infty$
(pertanyaan 6, dengan $b_n = n$) memberikan $\frac1n\sum \ln u_k \to
+\infty$, dan (G1) ($\eu^s \geq 1 + s$) mengirim rata-rata geometrinya
ke $+\infty$. Dengan $u_n = n$: $(n!)^{1/n} \to +\infty$.

\textbf{25.} (i) Kelengkapan masuk hanya lewat teorema limit
monoton, untuk menghasilkan limit pada pertanyaan 12 dan 20;
sedangkan teorema Cesàro--Stolz sendiri murni
pengelolaan $\varepsilon$, yang sah di atas $\Q$. (ii) Stolz mengganti
$\lim \frac{a_n}{b_n}$ dengan $\lim$ hasil bagi selisihnya,
persis seperti l'Hospital mengganti $\lim\frac fg$ dengan
$\lim\frac{f'}{g'}$ --- adapun kembaran diferensialnya bersandar pada teorema
nilai rata-rata pada \cref{ch:b1:derivative}. (iii) Heuristiknya: jika
$f(x) = x - a\,x^{p+1} + o(x^{p+1})$ pada titik tetap yang netral
$0$, maka $\frac{1}{u_{n+1}^p} - \frac{1}{u_n^p} \to pa$ dan
$u_n \sim (pan)^{-1/p}$: jadi kontak berorde $p + 1$ menghasilkan peluruhan
$n^{-1/p}$ --- sehingga makin rata grafiknya terhadap diagonalnya, makin
lambat jatuhnya. (iv) Konstantanya: (G2) memasok koefisien
kubiknya $\frac16$; lalu pemfaktoran pertanyaan 14 melipatduakannya
menjadi selisih teleskopnya $\frac13$; lalu Cesàro mengubah
$\frac{1}{u_n^2}$ menjadi $\frac n3$; dan membalikkannya lalu menarik akarnya
menghasilkan $\sqrt{3/n}$.

\end{solution}
